\section{Experimental Setup}
\label{sec:exp}

\subsection{Dataset and metrics}
We train and evaluate on COCO 2017~\cite{lin2014coco}. Training uses \texttt{train2017} (118k images, 860k bounding boxes); evaluation uses \texttt{val2017} (5k images). We report the standard COCO AP family:
\begin{itemize}
    \item $\mathrm{AP}$: mean average precision averaged over IoU $\in \{0.50, 0.55, \ldots, 0.95\}$ and 80 categories.
    \item $\mathrm{AP}_S$, $\mathrm{AP}_M$, $\mathrm{AP}_L$: same, restricted to objects with area $< 32^2$, $32^2$--$96^2$, $> 96^2$ pixels respectively.
    \item $\mathrm{AR}_{100}$: average recall at 100 detections per image.
\end{itemize}
$\mathrm{AP}_S$ is the headline metric for this work; the other entries are reported for completeness and to verify that the proposed mechanism does not degrade the easier categories.

\subsection{Architecture and training schedule}
The base architecture is RT-DETR-R50 with \texttt{ResNet50-vd} backbone~\cite{he2016resnet}. We add a stride-4 feature level (P2, 256 channels) to the backbone outputs, expand the input projection list, and increase the number of decoder queries from 300 to 500. Multi-scale training samples input sizes from $\{480, 512, 544, 576, 608, 640, 672, 704, 736, 768, 800\}$ uniformly each iteration, with $640$ kept as the modal sample. The decoder uses 6 layers; the feedback module fires after layer 0.

We initialize from the public \texttt{rtdetr\_r50vd\_6x\_coco.pth} weights~\cite{lyu2024rtdetr} and finetune for $13$ effective epochs on COCO train2017 with AdamW~\cite{loshchilov2019adamw}. Per-group learning rates:
\begin{center}
\begin{tabular}{ll}
    backbone (group 0) & $5\times10^{-7}$ \\
    \texttt{decoder.feedback} (group 1) & $5\times10^{-5}$ \\
    encoder norm/bias (group 2) & $5\times10^{-5}$ \\
    decoder norm/bias (group 3) & $5\times10^{-5}$ \\
    catch-all (group 4) & $5\times10^{-5}$ \\
\end{tabular}
\end{center}
Weight decay is $10^{-4}$ in all groups except norms/biases. We use AMP (\texttt{torch.cuda.amp}) for both speed and memory.

The 13 epochs are achieved as four chained SLURM jobs (one initial 1-epoch link aborted by an AMP dtype bug, then three 4-epoch links). Each chain link uses \texttt{--tuning} rather than \texttt{--resume}, which loads the model weights but reinitializes the optimizer and scheduler at the start of each link. As a side effect, the first epoch of each link looks like a small fresh-start dip; the cumulative AP trajectory in Section~\ref{sec:results} reflects this.

\subsection{Ablation protocol}
The \emph{causal} contribution of the feedback module is computed by toggling its forward pass at inference, with everything else held constant:
\begin{enumerate}
    \item Train v2 once. Save the final checkpoint $\theta^*$.
    \item Run the standard COCO val2017 evaluation on $\theta^*$ with feedback enabled. Record $\mathrm{AP}_S^{(\text{ON})}$.
    \item Reload the same $\theta^*$. Set \texttt{feedback.disabled = True} (Listing~\ref{lst:toggle}). Run the same evaluation. Record $\mathrm{AP}_S^{(\text{OFF})}$.
    \item Compute $\Delta\mathrm{AP}_S = \mathrm{AP}_S^{(\text{ON})} - \mathrm{AP}_S^{(\text{OFF})}$.
\end{enumerate}
The two evals differ only in whether the feedback module's forward returns the input memory unchanged. There is no retraining, no weight modification, no change in resolution or augmentation. Any difference in $\mathrm{AP}_S$ is attributable to the feedback computation. The same protocol is run on the v1 checkpoint, allowing a head-to-head v1-vs-v2 comparison of the causal contribution.

We additionally report a higher-resolution evaluation at $800\!\times\!800$ (feedback ON only) as a robustness check: if our small-object hypothesis is correct, the gain should grow when the input resolution is raised.

\subsection{Hardware}
Training was performed on the Bocconi University HPC cluster, on an A100 80\,GB GPU running in a 4g.40gb MIG slice (40\,GB visible memory) on the \texttt{stud} partition with QOS limit of one concurrent job per user. Per-epoch wall time was 2h25--2h45 depending on node load, with the full 13-epoch schedule completing in ${\approx}32$\,h of compute spread over four chain-resubmissions. Inference latency benchmarks (Section~\ref{sec:results-latency}) are reported on the same hardware.
